\subsection{Second-Moment Tensor, Isotropic Patterns, and Structural Exclusion}
  \label{subsec:second-moment}

  The Gram formulation developed in the previous section associates to each
  neutral generator $s_i$ a unit vector $\vec{u}_i\in S^2$.
  The geometry of the generating set can therefore be analysed through the
  configuration of these vectors on the unit sphere.

  A convenient global descriptor of this configuration is provided by the
  second-moment tensor
  \begin{equation}
    M=\sum_{i=1}^{k} \vec{u}_i \vec{u}_i^{\,T}.
  \end{equation}

  This tensor measures the directional distribution of the generators.
  Two qualitatively different regimes can occur.

  \paragraph{Anisotropic configurations.}
    If the tensor $M$ possesses distinct eigenvalues, the set of generators is
    anisotropically distributed on the sphere.
    Such configurations typically arise when generators concentrate along
    preferred directions.
    In this regime the admissible structures are closely related to
    generalised dihedral families, which can be continuously deformed while
    preserving the neutrality condition.

  \paragraph{Isotropic configurations.}
    A particularly important situation occurs when the second-moment tensor
    is proportional to the identity,
    \begin{equation}
      M \propto I .
    \end{equation}
    In this case the generators are distributed isotropically on the sphere.
    Isotropic patterns are much more constrained, since the neutrality and
    Gram relations fix the relative angles between generators up to discrete
    choices.

    Among isotropic configurations, a special role is played by orthogonal
    patterns in which the vectors $\vec{u}_i$ form mutually orthogonal
    directions.
    Such patterns arise naturally in the quaternion case discussed earlier
    and correspond to highly symmetric arrangements of generators.

    The isotropy condition therefore acts as a strong structural filter.
    Many candidate groups fail to satisfy the combined constraints imposed by
    neutrality, Gram rigidity, and isotropy.
    In particular, certain finite subgroups that might appear plausible from
    a purely group-theoretic perspective cannot realise a compatible
    generator geometry.

    A notable example is provided by the binary tetrahedral group $2T$.
    Although $2T$ is a finite subgroup of $SU(2)$, its character relations do
    not admit an isotropic configuration of neutral generators compatible
    with the Gram constraints.
    The corresponding second-moment tensor necessarily departs from
    isotropy, leading to an inconsistency with the admissibility conditions.

    The second-moment tensor thus provides a simple and effective diagnostic
    tool.
    By distinguishing isotropic from anisotropic generator configurations,
    it allows one to eliminate certain candidate groups and prepares the
    classification result presented in the next section.
